\documentclass{article}

\usepackage{neurips_2026}

\usepackage[utf8]{inputenc}
\usepackage[T1]{fontenc}
\usepackage{hyperref}
\usepackage{url}
\usepackage{booktabs}
\usepackage{amsfonts}
\usepackage{nicefrac}
\usepackage{microtype}
\usepackage{xcolor}

\title{Early Intent Classification for Humanoid Robot Interaction:\\
A Systematic Comparison of Text Encoders and Classifier Heads}

\author{%
  Author 1 \\
  School of Information Science and Technology \\
  ShanghaiTech University \\
  \texttt{author1@shanghaitech.edu.cn}
  \and
  Author 2 \\
  School of Information Science and Technology \\
  ShanghaiTech University \\
  \texttt{author2@shanghaitech.edu.cn}
}

\begin{document}

\maketitle

\begin{abstract}
Humanoid robots must decide whether a user utterance requires a verbal reply
(\texttt{chat}) or a physical action (\texttt{motion\_query}) before the
user finishes speaking.
We study this binary intent task on a 10{,}000-utterance dataset and compare
four text encoders (TF-IDF, CLIP, Qwen3-Embedding-0.6B, MiniLM-L6-v2)
with two lightweight classifier heads (logistic regression and a small MLP).
On full sentences, all learned models reach 99.4--99.8\% test accuracy, but
hard boundary cases reveal a speed--robustness trade-off: TF-IDF+LogReg is
fastest ($\sim$1.6\,ms) while MiniLM+LogReg is most robust on ambiguous
inputs (90.0\% hard accuracy).
For prefix-based early prediction, TF-IDF+LogReg achieves 94.5\% accuracy at
25\% of the utterance and 98.0\% overall across all prefix lengths, with a
mean correct decision latency of 33\% of the sentence at confidence
threshold~0.7.
We recommend TF-IDF+LogReg for real-time deployment and MiniLM+MLP when
robustness on partial inputs matters.
\end{abstract}

\section{Introduction}

In spoken interaction with humanoid robots, not every user query should trigger
the same response modality.
Informational questions (\emph{``What is machine learning?''}) call for
text generation, whereas imperative or demonstrative requests
(\emph{``Wave your hand''}) should activate motion planning.
Misclassifying \texttt{chat} as \texttt{motion\_query} causes unnecessary or
embarrassing robot gestures (\emph{false motion}); the reverse misses motion
requests (\emph{false text}).

Streaming speech adds latency: waiting for a complete sentence delays motion
preparation.
We therefore evaluate both \textbf{full-sentence classification} and
\textbf{prefix-based early prediction}, where the model sees only the first
25--100\% of words.

Our contributions are:
\begin{itemize}
  \item A balanced 10{,}000-utterance dataset for \texttt{chat} vs.\
        \texttt{motion\_query}, with a 30-example hard test set of boundary
        cases.
  \item A controlled $4 \times 2$ comparison (encoder $\times$ head) reporting
        accuracy, training time, and end-to-end inference latency.
  \item Prefix evaluation simulating incremental speech, measuring accuracy at
        four prefix ratios and early-decision latency under confidence
        thresholds.
\end{itemize}

\section{Method}

\subsection{Task and Labels}
Each utterance is labeled \texttt{chat} (verbal response) or
\texttt{motion\_query} (physical action).
We use macro-F1, false-motion rate (FMR), and false-text rate (FTR) alongside
accuracy.

\subsection{Dataset}
The dataset contains 10{,}000 utterances (5{,}000 per class): 2{,}000
LLM-generated (DeepSeek) plus 8{,}000 template-expanded samples.
We split 80/10/10 into train/val/test (8{,}000 / 1{,}000 / 1{,}000) with
stratification (seed~20260518).
A separate \emph{hard test} set of 30 manually curated boundary utterances
(e.g.\ \emph{``Can you explain how to wave properly?''}) evaluates
robustness.

\subsection{Encoders and Classifier Heads}
We fix the encoder and vary only the classification head:
\begin{itemize}
  \item \textbf{TF-IDF}: bag-of-words with 1--2-gram features (max 50{,}000).
  \item \textbf{CLIP}: \texttt{openai/clip-vit-base-patch32} text encoder.
  \item \textbf{Qwen3}: \texttt{Qwen/Qwen3-Embedding-0.6B}.
  \item \textbf{MiniLM}: \texttt{sentence-transformers/all-MiniLM-L6-v2}.
\end{itemize}
Heads: (1)~\textbf{LogReg}---sklearn logistic regression on cached
embeddings or TF-IDF vectors; (2)~\textbf{MLP}---LayerNorm $\to$ Linear(256)
$\to$ GELU $\to$ Linear(2), trained for 30 epochs with Adam.

Embedding vectors are pre-computed once on the full 10k corpus and reused for
both heads of the same encoder.

\subsection{Prefix (Early) Evaluation}
From each test utterance we extract word-level prefixes at ratios
25\%, 50\%, 75\%, and 100\%, yielding 4{,}000 evaluation rows.
Models classify each prefix independently.
We also measure \emph{early decision}: for confidence threshold~$\tau$,
the model commits at the earliest prefix whose predicted confidence exceeds
$\tau$, and we record the prefix ratio at which a \emph{correct} decision is
first made (decision latency).

\subsection{Baselines}
A keyword rule baseline matches motion-related verbs and nouns; it requires no
training and serves as a sanity check.

\section{Experiments}

All embedding experiments ran on an NVIDIA RTX~4060 Laptop GPU; TF-IDF ran on
CPU.
Latency is end-to-end per utterance (encode + classify), averaged over 200
test samples after 5 warmup runs.

\subsection{Full-Sentence Classification}

Table~\ref{tab:full-acc} shows test accuracy on the 1{,}000-example test split.
All learned models exceed 99.4\%; differences are small on in-distribution
data.

\begin{table}[h]
  \centering
  \caption{Test accuracy ($4 \times 2$ matrix) on full sentences.}
  \label{tab:full-acc}
  \begin{tabular}{lcc}
    \toprule
    Encoder & LogReg & MLP \\
    \midrule
    TF-IDF  & 0.994 & 0.997 \\
    CLIP    & 0.998 & 0.998 \\
    Qwen3   & 0.998 & 0.996 \\
    MiniLM  & 0.994 & 0.998 \\
    \bottomrule
  \end{tabular}
\end{table}

On the 30-example hard set (Table~\ref{tab:hard}), larger encoders do not
help: Qwen+MLP drops to 70.0\%, while MiniLM+LogReg reaches 90.0\% with zero
false-text errors.
TF-IDF+LogReg balances speed and hard-set performance (86.7\%).

\begin{table}[h]
  \centering
  \caption{Hard-set accuracy and error rates (30 boundary utterances).}
  \label{tab:hard}
  \small
  \begin{tabular}{lcccc}
    \toprule
    Model & Acc & FMR & FTR \\
    \midrule
    Keyword baseline & 0.967 & 0.000 & 0.033 \\
    TF-IDF+LogReg    & 0.867 & 0.033 & 0.100 \\
    TF-IDF+MLP       & 0.833 & 0.067 & 0.100 \\
    CLIP+LogReg      & 0.767 & 0.200 & 0.033 \\
    Qwen3+LogReg     & 0.767 & 0.200 & 0.033 \\
    Qwen3+MLP        & 0.700 & 0.133 & 0.167 \\
    MiniLM+LogReg    & \textbf{0.900} & 0.100 & \textbf{0.000} \\
    MiniLM+MLP       & 0.867 & 0.067 & 0.067 \\
    \bottomrule
  \end{tabular}
\end{table}

\subsection{Training Time and Inference Latency}

Table~\ref{tab:latency} summarizes the speed--accuracy trade-off.
TF-IDF+LogReg trains in 4.8\,s and infers in 1.6\,ms mean ($\sim$80$\times$
faster than Qwen3).
Qwen3 offers the highest embedding quality but 130\,ms latency makes it
unsuitable for streaming prefix prediction on resource-constrained robots.

\begin{table}[h]
  \centering
  \caption{Training time and single-utterance inference latency (ms).}
  \label{tab:latency}
  \small
  \begin{tabular}{lrrr}
    \toprule
    Model & Train & Mean & p95 \\
    \midrule
    TF-IDF+LogReg  & 4.8\,s  & 1.6  & 2.3  \\
    TF-IDF+MLP     & 25.7\,s & 1.7  & 2.8  \\
    MiniLM+MLP     & 30.4\,s & 15.1 & 19.7 \\
    CLIP+MLP       & 19.4\,s & 23.5 & 30.3 \\
    Qwen3+MLP      & 19.9\,s & 130.5 & 151.7 \\
    \bottomrule
  \end{tabular}
\end{table}

\subsection{Prefix-Based Early Prediction}

Table~\ref{tab:prefix} reports accuracy at each prefix ratio on 4{,}000 test
prefix rows (1{,}000 source utterances $\times$ 4 ratios).

\begin{table}[h]
  \centering
  \caption{Prefix accuracy by utterance completion ratio.}
  \label{tab:prefix}
  \begin{tabular}{lcccc}
    \toprule
    Model & 25\% & 50\% & 75\% & 100\% \\
    \midrule
    Keyword baseline & 0.829 & 0.927 & 0.915 & 0.917 \\
    TF-IDF+LogReg    & 0.945 & 0.987 & 0.995 & 0.994 \\
    Qwen3+MLP        & 0.943 & 0.938 & 0.990 & 0.996 \\
    MiniLM+MLP       & 0.777 & 0.969 & 0.993 & 0.998 \\
    \bottomrule
  \end{tabular}
\end{table}

TF-IDF+LogReg maintains 94.5\% accuracy after hearing only one quarter of an
utterance and reaches 98.0\% averaged over all prefix lengths.
MiniLM+MLP is fragile at 25\% (77.7\%, FMR 18.8\%) but recovers to 99.8\% at
full length---embedding models over-commit to motion from short, ambiguous
prefixes.
Qwen3+MLP matches TF-IDF at 25\% (94.3\%) but its 130\,ms inference cost
makes it impractical for streaming deployment despite comparable prefix
accuracy.

At confidence threshold $\tau = 0.7$ (Table~\ref{tab:early}), TF-IDF+LogReg
commits correctly after hearing a median of 30\% of words, with FMR~0.0\% and
FTR~2.2\%.
The keyword baseline decides equally early but misses 17.1\% of motion
requests (high FTR).

\begin{table}[h]
  \centering
  \caption{Early decision at confidence threshold $\tau = 0.7$ (1{,}000 sources).}
  \label{tab:early}
  \small
  \begin{tabular}{lrrrr}
    \toprule
    Model & Med.\ latency & Uncertain & FMR & FTR \\
    \midrule
    Keyword baseline & 0.30 & 0.0\% & 0.0\% & 17.1\% \\
    TF-IDF+LogReg    & 0.30 & 0.7\% & 0.0\% & 2.2\% \\
    Qwen3+MLP        & 0.30 & 0.0\% & 0.1\% & 5.4\% \\
    MiniLM+MLP       & 0.33 & 0.0\% & 16.5\% & 3.5\% \\
    \bottomrule
  \end{tabular}
\end{table}

\section{Conclusion}

We systematically compared four text encoders and two classifier heads for
humanoid intent routing on 10{,}000 utterances.
Full-sentence accuracy is high across all models, but hard boundary cases and
prefix evaluation expose meaningful differences.
\textbf{TF-IDF+LogReg} is the best default for real-time systems: 99.4\% test
accuracy, 1.6\,ms latency, and strong prefix performance (94.5\% at 25\%
completion).
\textbf{MiniLM+LogReg/MLP} is preferable when robustness on ambiguous or
partial inputs matters, at $\sim$15\,ms latency.
Large encoders (Qwen3, CLIP) offer marginal full-sentence gains but degrade on
hard cases and are too slow for streaming use.

Limitations include template-heavy training data (8{,}000 of 10{,}000 rows),
English-only coverage, and prefix simulation by word truncation rather than
real ASR output.
Future work will integrate streaming speech recognition and deploy the chosen
model on a physical humanoid platform with confidence-gated motion preparation.

\end{document}
